\documentclass[oneside,a4paper,14pt]{extarticle}
\usepackage[a4paper,letterpaper,top=20mm,bottom=20mm,left=20mm,right=10mm]{geometry}
\usepackage[russian]{babel}
\usepackage{indentfirst}
\usepackage{graphicx}
\usepackage{caption}
\usepackage{titlesec}
\usepackage{minted, fancyvrb}
\usepackage{hyperref}
\usepackage{enumitem}
\usepackage{amsmath}
\usepackage{fancyhdr}
\usepackage{colortbl}
\usepackage[dvipsnames]{xcolor}
\usepackage{tikz -timing}
\usepackage{tikz}
\usetikzlibrary{karnaugh}

\definecolor{LogisimKMapColor0}{RGB}{128,0,0}
\definecolor{LogisimKMapColor1}{RGB}{230,25,75}
\definecolor{LogisimKMapColor2}{RGB}{250,190,190}
\definecolor{LogisimKMapColor3}{RGB}{170,110,40}
\definecolor{LogisimKMapColor4}{RGB}{245,130,48}
\definecolor{LogisimKMapColor5}{RGB}{255,215,180}
\definecolor{LogisimKMapColor6}{RGB}{128,128,0}
\definecolor{LogisimKMapColor7}{RGB}{255,255,25}
\definecolor{LogisimKMapColor8}{RGB}{210,245,60}
\definecolor{LogisimKMapColor9}{RGB}{0,0,128}
\definecolor{LogisimKMapColor10}{RGB}{145,30,180}
\definecolor{LogisimKMapColor11}{RGB}{60,180,175}
\definecolor{LogisimKMapColor12}{RGB}{0,130,203}
\definecolor{LogisimKMapColor13}{RGB}{230,190,255}
\definecolor{LogisimKMapColor14}{RGB}{170,255,195}
\definecolor{LogisimKMapColor15}{RGB}{240,50,230}

% Форматирование листингов с кодом
\setminted{style = rainbow_dash, fontsize = \small} % https://pygments.org/styles/

% Форматирование заголовков
\titleformat{\section}{\normalsize\bfseries}{\thesection}{1em}{}
\titleformat{\subsection}{\normalsize\bfseries}{\thesubsection}{1em}{}
\titleformat{\subsubsection}{\normalsize\bfseries}{\thesubsubsection}{1em}{}

% Интерлиньяж и абзац
\renewcommand\baselinestretch{1.33}

\setlength{\parindent}{1.25cm}  % длина красной строки

% Для всех списков
\setlist[enumerate]{
  left=0.5\parindent,       % отступ слева
  label=\arabic*.,       % цифры
  itemsep=0pt,           % расстояние между пунктами
  topsep=5pt,            % отступ сверху
  partopsep=0pt,         % дополнительный отступ сверху, если абзац до списка
  parsep=0pt             % отступ между абзацами внутри пункта
}

\setlist[itemize]{
  left=0.5\parindent,       % отступ слева
  itemsep=0pt,           % расстояние между пунктами
  topsep=5pt,            % отступ сверху
  partopsep=0pt,
  parsep=0pt
}

% Гиперссылки
\hypersetup{
  colorlinks=true,
  linkcolor=black,
  urlcolor=blue,
  pdfborder={0 0 0},
  pdftitle={Синтез комбинационной схемы},
  pdfauthor={Черкасов А.А. и Зинин В.А.},
}

\begin{document}

\newpage
\thispagestyle{empty}
\begin{center}
  МИНИСТЕРСТВО НАУКИ И ВЫСШЕГО ОБРАЗОВАНИЯ РОССИЙСКОЙ ФЕДЕРАЦИИ ФЕДЕРАЛЬНОЕ ГОСУДАРСТВЕННОЕ БЮДЖЕТНОЕ УЧРЕЖДЕНИЕ ВЫСШЕГО ОБРАЗОВАНИЯ\\
  <<ВЯТСКИЙ ГОСУДАРСТВЕННЫЙ УНИВЕРСИТЕТ>>\\
  Институт математики и информационных систем\\
  Факультет автоматики и вычислительной техники\\
  Кафедра электронных вычислительных машин
\end{center}
\vspace{10mm}

\hfill
\begin{tabular}{l}
  \footnotesize Дата сдачи на проверку:                                          \\
  \footnotesize <<\rule[-1mm]{5mm}{0.10mm}\/>>\rule[-1mm]{20mm}{0.10mm}\ 2025 г. \\
  \footnotesize Проверено:                                                       \\
  \footnotesize <<\rule[-1mm]{5mm}{0.10mm}\/>>\rule[-1mm]{20mm}{0.10mm}\ 2025 г. \\
\end{tabular}
\vfill

\begin{center}
  Отчёт по лабораторной работе №1\\
  по дисциплине\\
  <<Цифровая схемотехника>>\\
\end{center}
\vspace{25mm}
\noindent
\begin{tabular}{ll}
  Разработали студенты гр. ИВТб-2301-05-00 & \hspace{5mm}\rule[-1mm]{30mm}{0.10mm}\,/Черкасов А. А./ \\
                                           & \hspace{5mm}\rule[-1mm]{30mm}{0.10mm}\,/Зинин В. А./    \\
                                           & \hspace{12.5mm}\footnotesize(подпись)                   \\
  Преподователь                            & \hspace{5mm}\rule[-1mm]{30mm}{0.10mm}\,/Мельцов В. Ю./  \\
                                           & \hspace{12.5mm}\footnotesize(подпись)                   \\
\end{tabular}

\noindent
\begin{tabular}{lp{58mm}r}
  Работа защищена &  & \hspace{5mm}<<\rule[-1mm]{5mm}{0.10mm}\/>>\rule[-1mm]{30mm}{0.10mm}\ 2026 г.
\end{tabular}
\vfill

\begin{center}
  Киров\\
  2026
\end{center}

\newpage\thispagestyle{plain}

\section{Цель лабораторной работы}
Получить навыки синтеза комбинационных схем в определённом логическом базисе. Синтезировать логическую схему по заданной таблице истинности, используя базис <<И, ИЛИ, НЕ>>.

\section{Таблица истинности}
\begin{center}
  Таблица 1 --- Таблица истинности \\
  \vspace{5mm}
  \begin{tabular}{cccc|c}
    $x_0$ & $x_1$ & $x_2$ & $x_3$ & $F$ \\
    \hline
    $0$   & $0$   & $0$   & $0$   & $0$ \\
    $0$   & $0$   & $0$   & $1$   & $0$ \\
    $0$   & $0$   & $1$   & $0$   & $0$ \\
    $0$   & $0$   & $1$   & $1$   & $0$ \\
    $0$   & $1$   & $0$   & $0$   & $0$ \\
    $0$   & $1$   & $0$   & $1$   & $1$ \\
    $0$   & $1$   & $1$   & $0$   & $0$ \\
    $0$   & $1$   & $1$   & $1$   & $0$ \\
    $1$   & $0$   & $0$   & $0$   & $0$ \\
    $1$   & $0$   & $0$   & $1$   & $0$ \\
    $1$   & $0$   & $1$   & $0$   & $0$ \\
    $1$   & $0$   & $1$   & $1$   & $0$ \\
    $1$   & $1$   & $0$   & $0$   & $0$ \\
    $1$   & $1$   & $0$   & $1$   & $0$ \\
    $1$   & $1$   & $1$   & $0$   & $1$ \\
    $1$   & $1$   & $1$   & $1$   & $1$ \\
  \end{tabular}
\end{center}

\section{Уравнение функции $F$}
\begin{equation*}
F = (\overline{x_0} \cdot x_1 \cdot \overline{x_2} \cdot x_3) + (x_0 \cdot x_1 \cdot x_2 \cdot \overline{x_3}) + (x_0 \cdot x_1 \cdot x_2 \cdot x_3)
\end{equation*}

\section{Построим карту Карно для минимизации КУ}

\begin{center}
  Таблица 2 --- Карта Карно \\
  \vspace{5mm}
  \begin{tikzpicture}[karnaugh,disable bars,x=1\kmunitlength,y=1\kmunitlength,kmbar left sep=1\kmunitlength,grp/.style n args={4}{#1,fill=#1!30,minimum width= #2\kmunitlength,minimum height=#3\kmunitlength,rounded corners=0.2\kmunitlength,fill opacity=0.6,rectangle,draw}]
    \karnaughmap{4}{$F$}{{$x_0$}{$x_2$}{$x_1$}{$x_3$}}
    {0001000000000011}{
      \draw[kmbox] (-0.5,4.5)
      node[below left]{$x_0$, $x_1$}
      node[above right]{$x_2$, $x_3$} +(-0.2,0.2)
      node[above left]{$F$};\draw (0,4) -- (-0.7,4.7);
      \foreach \x/\1 in %
        {0/00,1/01,2/11,3/10} {
          \node at (\x+0.5,4.2) {\1};
        }
      \foreach \y/\1 in %
        {0/00,1/01,2/11,3/10} {
          \node at (-0.4,-0.5-\y+4) {\1};
        }
      \node[grp={LogisimKMapColor0}{0.8}{0.8}](n0) at(1.5,2.5) {};
      \node[grp={LogisimKMapColor1}{1.8}{0.8}](n1) at(3,1.5) {};
    }
  \end{tikzpicture}
\end{center}
Получаем минимальное выражение для функции $F$:
\begin{equation*}
  F =  \overline{x_0}  \cdot x_1 \cdot  \overline{x_2}  \cdot x_3+x_0 \cdot x_1 \cdot x_2
\end{equation*}

\section{Построим функциональную схему в базисе <<И, ИЛИ, НЕ>>}
\begin{figure}[H]
  \centering
  \includegraphics[width=0.9\textwidth]{pics/functional_diagram.png}
  \caption*{Рисунок 1 --- Функциональная схема устройства в базисе <<И, ИЛИ, НЕ>>}
\end{figure}
\newpage
\section{Построим принципиальную схему устройства}

\begin{figure}[H]
  \centering
  \includegraphics[width=0.9\textwidth]{pics/principle_diagram.png}
  \caption*{Рисунок 2 --- Принципиальная схема устройства в базисе <<И, ИЛИ, НЕ>>}
\end{figure}
\newpage

Определим контакты питания для микросхем ТТЛ:
\begin{center}
  Таблица 4 --- Таблица питания для микросхем ТТЛ \\
  \vspace{5mm}
  \begin{tabular}{|c|c|c|}
    \hline
    Микросхема     & +5V & 0V \\
    \hline
    D1, D2, D3, D4 & 14 & 07   \\
    \hline
  \end{tabular}
\end{center}
\section*{Вывод}

В данной лабораторной работе мы получили навыки синтеза комбинационных схем в базисе <<И, ИЛИ, НЕ>>. Мы синтезировали логическую схему по заданной таблице истинности, используя карту Карно для минимизации функции. В результате мы получили минимальное выражение для функции $F$ и построили функциональную и принципиальную схемы устройства.


\end{document}
